\begin{problem}[Ideal-sheaf quotient]\label{prob:vi15-p16}
If $\mathcal I\subseteq\mathcal O$ is a sheaf of ideals, prove $\mathcal O/\mathcal I$ is a sheaf of rings and $(\mathcal O/\mathcal I)_x\cong\mathcal O_x/\mathcal I_x$.
\end{problem}

\ifdefined\FullProblemDossiers
\paragraph{Why this problem matters.}
This is the algebraic mechanism behind closed subschemes.

\paragraph{Useful ideas.}
Translate the sheaf statement to local representatives or stalks whenever possible.  Keep track of whether the construction is sectionwise or requires sheafification.

\paragraph{Guided hints.}
Check multiplication on local representatives and use quotient stalks.

\begin{solution}
For every open $U$, $\mathcal O(U)/\mathcal I(U)$ is a ring because $\mathcal I(U)$ is an ideal. Restriction maps are ring homomorphisms and preserve $\mathcal I$, so the sectionwise quotients form a presheaf of rings. Its sheafification is the quotient sheaf $\mathcal O/\mathcal I$.

The ring operations on the sheafification are defined locally: if sections are represented on a cover by classes $[a_i]$ and $[b_i]$, define their sum and product by $[a_i+b_i]$ and $[a_ib_i]$. Changing representatives by elements of $\mathcal I$ changes these expressions by elements of $\mathcal I$, so the operations are well defined and glue. Hence $\mathcal O/\mathcal I$ is a sheaf of rings.

Finally,
\[
(\mathcal O/\mathcal I)_x\cong\mathcal O_x/\mathcal I_x
\]
by the quotient-stalk theorem, and the isomorphism respects multiplication because it sends the local class of $a$ to the class of the germ $a_x$.
\end{solution}

\paragraph{What happened mathematically.}
This is the algebraic mechanism behind closed subschemes.

\paragraph{Extensions.}
Formulate the same statement for sheaves of modules.  Where meaningful, test the statement on a discrete space and then on the exponential-map example to see the difference between sectionwise and local behavior.

\paragraph{Related problems.}
Compare VI.15.P1 and VI.15.P2, the two preserved corpus problems.  VI/16 will reorganize these constructions into general exact-sequence language.

\paragraph{Provenance.}
New synthesis; previews VI/22.
\else
\paragraph{Provenance.} New synthesis; previews VI/22.
\fi
